% PUV Pipeline — L1 and L2 Methods Section
% Drop into your paper's methods section. Adjust \subsection depth as needed.
% Requires: amsmath, siunitx (optional, for units)

\subsection{Nearshore wave measurements}

Nearshore wave and current dynamics were measured using bottom-mounted
Nortek Vector acoustic Doppler velocimeters (ADVs) configured in
pressure--velocity (PUV) mode, sampling continuously at \SI{2}{\hertz}.
Instruments were deployed along MOP (Monitoring and Prediction) transect
lines at nominal depths of \SIrange{5}{7}{\meter}, with two cross-shore
positions per transect. Deployment metadata---including instrument heading,
clock drift, and pressure sensor height above the bed ($z_s$, typically
\SIrange{0.75}{0.77}{\meter})---were obtained from field deployment notes.

\subsection{Level 1: Quality control and coordinate rotation}

Raw burst files (.dat/.sen/.hdr) were merged into continuous time series
at the \SI{2}{\hertz} sampling rate. Clock drift corrections were applied
linearly over each deployment based on measured drift at recovery. A
battery cutoff detection algorithm identified intermittent sampling
caused by battery depletion by flagging gaps exceeding \SI{1}{\second}
between consecutive sensor records; all data following the first such
gap were discarded.

Quality control filters removed samples where: (1) pitch or roll exceeded
$\pm\SI{5}{\degree}$; (2) pressure fell outside the range
$[\overline{P}/2,\; 2\overline{P}]$, where $\overline{P}$ is the deployment median
pressure; or (3) the minimum beam correlation dropped below 70\%, following
Nortek recommendations. Flagged samples were set to NaN.

A variability-based tilt quality control was applied rather than a fixed
angular threshold. The rolling standard deviation of pitch and roll was
computed over 60-second windows; samples where the variability exceeded
\SI{2}{\degree} were flagged as unstable. An absolute tilt cap of
\SI{30}{\degree} was retained for severely compromised instruments.
Samples with stable but non-zero tilt (e.g., from pipes bent during storm
events) were retained and corrected via a sample-by-sample
three-dimensional rotation using the measured pitch ($\phi$) and roll
($\theta$) angles:
%
\begin{equation}
  \mathbf{v}_\text{corrected} =
  \mathbf{R}_\text{roll}(\theta)\;\mathbf{R}_\text{pitch}(\phi)\;
  \mathbf{v}_\text{raw},
\end{equation}
%
where
\[
  \mathbf{R}_\text{roll} = \begin{bmatrix}
    1 & 0 & 0 \\
    0 & \cos\theta & -\sin\theta \\
    0 & \sin\theta &  \cos\theta
  \end{bmatrix}, \qquad
  \mathbf{R}_\text{pitch} = \begin{bmatrix}
    \cos\phi & 0 & \sin\phi \\
    0 & 1 & 0 \\
    -\sin\phi & 0 & \cos\phi
  \end{bmatrix}.
\]
This tilt correction was applied before the heading rotation to buoy
coordinates, ensuring that the velocity vector was first returned to the
local vertical frame before projection onto the geographic coordinate
system.

Velocity components were rotated from the instrument XYZ frame to a
geographic buoy coordinate system ($+x$ west, $+y$ north, $+z$ up;
left-handed, consistent with the MOP convention) using the magnetic
heading from deployment notes and the magnetic declination computed via
the International Geomagnetic Reference Field (IGRF-13) model at the
deployment location and epoch.

\subsection{Level 2: Spectral analysis}

\subsubsection{Segmentation and preprocessing}

The continuous L1 time series were divided into non-overlapping
\SI{1}{\hour} segments of 7200 samples at \SI{2}{\hertz}. The 1-hour
segment length matches the cadence of the regional CDIP MOP wave
hindcast and gives the longer averaging window that segment-mean
velocity diagnostics benefit from (random Doppler error in
ū decreases as $1/\sqrt{N}$, so the 7200-sample segment-mean noise
floor is approximately $1.87\times$ smaller than the legacy
2048-sample 17-minute alternative). A within-hour stationarity audit
on the existing record confirmed that bulk wave parameters
($H_s$, $T_p$, $T_{m02}$) are stable on the 1-hour scale at the
sites and depths considered here. Segments in which more than 10\%
of pressure or velocity samples were flagged as NaN were discarded;
for segments with $\leq$10\% NaN, gaps were filled by linear
interpolation before further processing. Prior to spectral analysis,
velocities were rotated to a shore-normal coordinate frame
($+x$ onshore, $+y$ alongshore north) using the shore-normal angle
retrieved from the CDIP THREDDS server for the corresponding MOP station.

For each valid segment, the mean pressure was used to compute the total
water column depth $H = \overline{P} \times 10^4 / (\rho g) + z_s$, where
$\overline{P}$ is the segment-mean pressure in dBar, $\rho =
\SI{1025}{\kilogram\per\meter\cubed}$, $g = \SI{9.81}{\meter\per
\second\squared}$, and $z_s$ is the pressure sensor height above the bed.
Mean cross-shore, alongshore, and vertical velocities and temperature
were recorded for each segment prior to detrending. The pressure time
series was converted from dBar to meters of water head, and
all channels (pressure, $u$, $v$, $w$) were linearly detrended.

\subsubsection{Spectral estimation}

Power spectral densities were estimated using the multi-taper method
of \citet{thomson1982} with time-bandwidth product $NW = 4$, giving
$K = 7$ orthogonal Slepian tapers and approximately
$2K = 14$ degrees of freedom per frequency bin. The
1-hour segment provides a native frequency resolution of
$\Delta f \approx \SI{2.78e-4}{\hertz}$ across the full
0--\SI{1}{\hertz} range. Cross-spectra between pressure and the two
horizontal velocity components were computed using the same tapers
for directional analysis. The multi-taper choice was made for its
optimal bias--variance trade-off relative to the legacy Welch
estimator: at constant total record length, MTM with $NW=4$ on
the 1-hour record gives an effective bandwidth and bias comparable
to Welch with 256-sample sub-segments while preserving fine
frequency structure that Welch averaging smears \citep{percival1993}.

\subsubsection{Pressure correction}

The measured pressure spectrum was converted to a surface elevation
spectrum using the linear wave theory transfer function:
%
\begin{equation}
  K_p(f) = \frac{\cosh(k z_s)}{\cosh(k H)}
\end{equation}
%
where $k(f)$ is the wavenumber obtained by solving the linear dispersion
relation $\omega^2 = g k \tanh(kH)$ via Newton's method. The surface
elevation spectrum was computed as $S_\eta(f) = S_p(f) / K_p^2(f)$.
Frequency bins where $K_p < 0.1$ were set to zero to suppress
noise amplification at high frequencies; the resulting cutoff frequency
$f_\text{cut}$ was recorded for each segment.

\subsubsection{Bulk wave parameters}

Significant wave height was computed as $H_s = 4\sqrt{m_0}$, where
$m_0 = \int S_\eta(f)\,df$ is the zeroth spectral moment, integrated
separately over the sea--swell band (\SIrange{0.04}{0.25}{\hertz})
and infragravity band (\SIrange{0.004}{0.04}{\hertz}). Peak period
$T_p = 1/f_p$ was determined from the frequency of maximum $S_\eta$
within the sea--swell band, and the mean zero-crossing period was
computed as $T_{m02} = \sqrt{m_0/m_2}$. The energy-weighted mean wave
direction was obtained from the first-harmonic directional coefficients
$a_1$ and $b_1$ derived from the pressure--velocity cross-spectra.
Total energy flux was computed as
$F = \rho g \int c_g(f)\, S_\eta(f)\, df$,
where $c_g$ is the group velocity from linear wave theory.

\subsubsection{Near-bed velocity and stress}

Near-bed orbital velocities were estimated by scaling the measured
velocity spectrum at sensor height to the bed using linear wave theory.
Each frequency bin of the velocity FFT was multiplied by
$\cosh(kH) / \cosh[k(H - z_s)]$, and the result was inverse-transformed
to obtain the near-bed velocity time series \citep[e.g.,][]{guza1980}.
The RMS near-bed orbital velocity $U_b$ and wave orbital excursion $A_w$
were computed from the bed velocity time series for each segment.

Wave-induced bed shear stress was estimated using the Swart (1974)
friction factor formulation with a bed roughness of $k_s = 2.5 D_{50}$,
where $D_{50} = \SI{0.25}{\milli\meter}$ (to be updated with site-specific
grain size distributions from laser particle analysis). Reynolds stress
components ($\overline{u'w'}$, $\overline{v'w'}$, $\overline{u'v'}$)
and turbulent kinetic energy were computed from the detrended velocity
fluctuations. Velocity moment statistics including skewness $S_k$ and
velocity asymmetry $A_s$ (Hilbert form, negative for the forward-pitched
shoaling wave shape) were computed following
\citet{elgar1985,hoefel2003} for characterizing wave nonlinearity
relevant to sediment transport.

\subsection{Validation against MOP wave model}

PUV-derived wave parameters were validated against hourly wave data from
the CDIP Monitoring and Prediction (MOP) system, a spectral wave model
initialized by offshore CDIP buoy observations. The MOP model provides
wave energy density spectra $S_\eta(f)$ at \SI{10}{\meter} depth along
each MOP transect line. MOP spectra were shoaled to the PUV instrument
depth using linear wave theory energy conservation:
%
\begin{equation}
  S_\eta(f, h_\text{PUV}) = S_\eta(f, 10\,\text{m}) \times
  \frac{c_g(f, 10\,\text{m})}{c_g(f, h_\text{PUV})}
\end{equation}
%
Shoaled MOP wave heights were computed by integrating the shoaled spectrum
over the sea--swell band using the (non-uniform) MOP frequency bandwidth
vector. PUV segment values (\SI{17}{\minute} resolution) were interpolated
to the MOP hourly time grid for comparison. Agreement was quantified using
$R^2$, RMSE, bias (PUV $-$ MOP), and the Willmott (1981) skill score.

Frequency-dependent validation was performed by computing the mean
spectral ratio $S_\eta^\text{PUV}(f) / S_\eta^\text{MOP,shoaled}(f)$
across all matched hours, binned into low-swell
(\SIrange{0.04}{0.08}{\hertz}), mid-swell (\SIrange{0.08}{0.14}{\hertz}),
and high-swell (\SIrange{0.14}{0.25}{\hertz}) bands. This diagnostic
isolates whether systematic biases originate from the pressure transfer
function (which amplifies high-frequency energy most aggressively) or
from other sources.

To investigate the origin of systematic biases, four candidate
mechanisms were tested: (1)~nonlinear shoaling via triad interactions,
assessed by correlating the spectral ratio with the Ursell number
$\mathit{Ur} = H_s L_p^2 / h^3$; (2)~directional narrowing during
shoaling, assessed by comparing the first-harmonic directional spread
$\sigma_1 = \sqrt{2(1 - \sqrt{a_1^2 + b_1^2})}$ between PUV and MOP;
(3)~bound long wave contamination, tested by examining the scaling of
low-frequency excess energy with $H_s^2$ and the frequency structure of
the excess relative to the expected group frequency; and (4)~MOP model
spectral peak broadening, quantified by comparing the Goda peakedness
parameter $Q_p = (2/m_0^2) \int f\, S_\eta^2(f)\, df$ and the
half-power spectral bandwidth between PUV and shoaled MOP spectra.

\subsection{Level 3: Wave forcing characterization}

\subsubsection{Frequency-band energy decomposition}

The L2 surface elevation spectrum was integrated over three frequency
bands to characterize wave energy partitioning: infragravity
(0.004--0.04~Hz), swell (0.04--0.10~Hz), and sea (0.10--0.25~Hz).
Band-specific significant wave height, energy flux, and mean period
were computed for each 17-minute segment. A dominant band flag was
assigned based on which band carried the largest energy flux.

\subsubsection{Storm event detection}

Storm events were identified from the $H_s$ time series using a
threshold-exceedance method with configurable parameters (default:
$H_s > 1.5$~m sustained for $\geq$6~hours, events separated by
$<$12~hours merged). To characterize wave forcing during PUV data
gaps (e.g., from instrument damage during storms), hourly MOP wave
model data were loaded for a window extending 7~days beyond the PUV
record. Storm events detected in the MOP record but absent from the
PUV data were flagged as gap-period storms, providing continuous
forcing context even when the in-situ instruments were compromised.

\subsubsection{Transport proxies}

Bottom energy flux was computed as the spectral integral of group
velocity times surface elevation spectral density over the sea--swell
band. The Shields parameter $\theta = \tau_b / [(\rho_s - \rho) g D_{50}]$
was computed per segment, with the critical Shields threshold from the
Soulsby--Whitehouse formula. A mobilization flag indicated segments
exceeding the critical threshold. Cumulative bottom energy flux was
integrated over the deployment duration for comparison with
morphological change between beach surveys. The Rouse number
$P = w_s / (\kappa u_*)$ classified the dominant transport mode
(bedload vs.\ suspended load).

\subsubsection{Tidal current decomposition}

Tidal harmonic analysis was performed on the segment-mean velocities
using the \texttt{t\_tide} package, with complex velocity input
($u + iv$) enabling simultaneous analysis of cross-shore and
alongshore tidal constituents. The tidal depth signal was obtained
from NOAA tide predictions for the Scripps Pier gauge (station
9410230), which provided a smooth, gap-free tidal reference that
outperformed the \texttt{t\_tide} fit to the noisy PUV pressure
record ($R = 0.995$ vs.\ $R = 0.987$ for TBR23). PUV instrument
clocks were confirmed to be in UTC via cross-correlation with the
NOAA record (optimal lag $= 0.0$~hours). Subtidal residual currents
(observed minus tidal) captured wave-driven mean flows including
undertow and longshore currents.

\subsection{Mean cross-shore flow validation}
\label{sec:mean_flow_validation}

The segment-mean cross-shore velocity $\overline{u}$ was tested
against three competing hypotheses---real wave-driven Stokes
return flow, a constant per-instrument calibration offset, or
random Doppler noise---using a multi-test framework on every PUV
deployment in the catalog (38 instruments across Torrey Pines, Los
Pe\~nasquitos lagoon, Solana Beach, SIO Pier, Catalina Island, and
the legacy Ruby2D cross-shore array).

\subsubsection{Linear scaling against $H_s^2$}
For each instrument, the segment-mean cross-shore velocity was fit
as a linear function of significant wave height squared:
%
\begin{equation}
  \overline{u} = \alpha\, H_s^2 + \beta + \varepsilon,
  \label{eq:mean_flow_fit}
\end{equation}
%
restricted to segments with $H_s \geq \SI{0.2}{\meter}$. The slope
$\alpha$ and intercept $\beta$ provide the discriminating
coefficients. Stokes return-flow theory \citep{longuethiggins1953}
predicts a depth-averaged Eulerian compensation flow
$\overline{u} \approx -gH_s^2 / (16ch)$ for shallow water with
$c = \sqrt{gh}$, giving $\alpha_\text{theory} = -g/(16ch)$, which
ranges from $\approx -0.018$ at $h = \SI{5}{\meter}$ to
$\approx -0.003$ at $h = \SI{15}{\meter}$ in m/s per
m\textsuperscript{2}\,of $H_s^2$. The theoretical intercept is
$\beta_\text{theory} = 0$. The three hypotheses make different
predictions: real Stokes return flow gives $\alpha$ matching theory
in sign and within an order-1 factor with $\beta$ centered on zero
across many instruments; a constant per-instrument offset gives
$\alpha \approx 0$ with $\beta$ scattering at the offset magnitude;
random Doppler noise leaves both coefficients unconstrained but
shrinking toward zero with longer averaging.

\subsubsection{Three robustness checks}
Three follow-up fits address potential rebuttals to the
$H_s^2$-scaling argument. First, the same fit is run on the
alongshore component $\overline{v}$ to test whether $\alpha$ is
specific to cross-shore physics or reflects generic storm-coherent
forcing (wind stress, atmospheric pressure setup, alongshore
gradients) that also correlates with $H_s$. Second, the fit is
restricted to $H_s \geq \SI{1}{\meter}$, the wave-forcing-dominant
regime where Stokes theory applies most cleanly, and the change in
$R^2$ is reported. Third, an explicit wave-direction-discrimination
test extends the alongshore model to
%
\begin{equation}
  \overline{v} = \alpha_{v0}\, H_s^2 + \alpha_{v1}\, \sin(2\theta_\text{rel})\, H_s^2 + \beta_v,
  \label{eq:wave_dir_check}
\end{equation}
%
where $\theta_\text{rel}$ is the wave incidence angle relative to
shore normal (taken directly from the L2 mean wave direction, which
is computed in the rotated frame). A coordinate-rotation error
$\theta_\text{err}$ in the shore-normal angle predicts
$\alpha_{v0} = \alpha_u \tan(\theta_\text{err})$ with
$\alpha_{v1} \approx 0$, while a real radiation-stress-driven
longshore current predicts $\alpha_{v1}$ dominant with the sign
matching local wave climatology. The implied
$\theta_\text{err} = \arctan(\alpha_{v0}/\alpha_u)$ provides an
upper bound on any actual rotation error.

\subsubsection{Auxiliary diagnostics}
Three additional tests provide independent corroboration. (1)~Tidal
phase modulation of $\overline{u}$ was computed by binning
segment-mean velocities into 16 phase bins of band-passed depth
(\SIrange{0.0019}{0.0093}{\hertz}, $\sim$30~min to 3~hr period),
stratified by $H_s$. A constant instrumental offset has zero
modulation amplitude in any $H_s$ stratum. (2)~Cross-instrument
correlations of $\overline{u}(t)$ between independently-calibrated
Vectors at common timestamps test whether shared time-varying
variance exists; independently-calibrated instruments cannot
produce coincidentally biased records. (3)~The Reynolds-stress
component $\overline{u'w'}$, derived from the same record via an
independent processing path (segment-mean of fluctuating velocity
covariance rather than the segment mean itself), was correlated
with $\overline{u}$. Independent random noise in either path would
push the correlation toward zero.

\subsubsection{Segmentation independence}
The mean-flow validation is run at both the canonical 1-hour
segmentation and the legacy 17-minute segmentation. Random
instrument noise predicts that $\overline{u}$ estimates drift
toward zero as the averaging window grows; if the signal is real
wave-driven flow, $\alpha$ is approximately segmentation-
independent (kinematic relationship) while $|\beta|$ tightens
slightly toward zero. The segmentation comparison is therefore a
direct falsification test against the noise hypothesis.
